% x86-64 GAS Assembly Quick Reference Card — PRO
% Compile: pdflatex card_pro.tex
\documentclass[landscape,a4paper]{article}
\usepackage[landscape,margin=0.4cm,top=0.5cm,bottom=0.4cm]{geometry}
\usepackage{multicol}
\usepackage{booktabs}
\usepackage{array}
\usepackage{xcolor}
\usepackage{tcolorbox}
\usepackage{listings}
\usepackage[T1]{fontenc}
\usepackage{lmodern}

\pagestyle{empty}
\tcbuselibrary{skins}

%%% Font sizes %%%
\makeatletter
\renewcommand\normalsize{\@setfontsize\normalsize{8pt}{9.5pt}}
\renewcommand\small{\@setfontsize\small{7.5pt}{9pt}}
\renewcommand\footnotesize{\@setfontsize\footnotesize{7pt}{8.5pt}}
\renewcommand\scriptsize{\@setfontsize\scriptsize{6pt}{7pt}}
\renewcommand\tiny{\@setfontsize\tiny{5pt}{6pt}}
\makeatother

%%% Operand type colours %%%
\definecolor{tL}{RGB}{0,80,180}
\definecolor{tM}{RGB}{180,30,30}
\definecolor{tR}{RGB}{0,110,50}
\newcommand{\Lc}{\textcolor{tL}{\textbf{L}}}
\newcommand{\Mc}{\textcolor{tM}{\textbf{M}}}
\newcommand{\Rc}{\textcolor{tR}{\textbf{R}}}

%%% tcolorbox style %%%
\tcbset{
  mysec/.style={
    colback=white, colframe=black!60,
    boxrule=0.3pt, arc=1pt,
    fonttitle=\normalsize\bfseries,
    left=3pt, right=3pt, top=1pt, bottom=1pt,
    before skip=2pt, after skip=2pt,
  }
}

%%% Listing style %%%
\lstdefinestyle{asmcode}{
  basicstyle=\ttfamily\footnotesize,
  commentstyle=\color{gray}\itshape,
  morecomment=[l]{\#},
  showstringspaces=false,
  breaklines=true,
  frame=none,
  xleftmargin=0pt, xrightmargin=0pt,
  aboveskip=0pt, belowskip=0pt,
  columns=fullflexible,
  keepspaces=true,
  tabsize=4,
}

\begin{document}
\small
\setlength{\tabcolsep}{2.5pt}
\setlength{\extrarowheight}{1pt}
\setlength{\columnsep}{4pt}
\setlength{\columnseprule}{0.2pt}
\renewcommand{\arraystretch}{1.0}

% ====================== TITLE ======================
\noindent
{\large\bfseries x86-64 GAS Assembly Quick Reference Card\enspace\texttt{PRO}}
\hfill
{\footnotesize
  \Lc\ = Immediate/Literal \quad
  \Mc\ = Memory \quad
  \Rc\ = Register \quad
  \{S\} = Size suffix: b(byte) w(word) l(dword) q(qword)}

\vspace{2pt}

% ===================== PART 1 =====================
\begin{tcolorbox}[mysec,
  title={Part 1: Instruction Set --- AT\&T: \texttt{op src, dst}
         \quad$|$\quad Intel: \texttt{op dst, src}}]

\begin{multicols}{2}
\setlength{\parskip}{0pt}
\setlength{\parsep}{0pt}

% ---------- AT&T Syntax ----------
{\small\bfseries\sffamily AT\&T Syntax}\par\vspace{-3pt}
{
\begin{tabular}{@{}l l l l p{4.2cm}@{}}
\toprule
\textbf{Mnem.} & \textbf{Op1} & \textbf{Op2} & \textbf{Op3}
  & \textbf{Description} \\
\midrule
\texttt{mov\{S\}}  & \Rc/\Mc/\Lc & \Rc/\Mc & --- & Copy src$\to$dst \\
\texttt{movsx}     & \Rc/\Mc     & \Rc     & --- & Move with sign-extension \\
\texttt{movzx}     & \Rc/\Mc     & \Rc     & --- & Move with zero-extension \\
\texttt{lea}       & \Mc         & \Rc     & --- & Load effective address \\
\texttt{xchg\{S\}} & \Rc/\Mc     & \Rc/\Mc & --- & Exchange two operands \\
\midrule
\texttt{add\{S\}}  & \Rc/\Mc/\Lc & \Rc/\Mc & ---
  & $dst \mathrel{+}= src$ \\
\texttt{sub\{S\}}  & \Rc/\Mc/\Lc & \Rc/\Mc & ---
  & $dst \mathrel{-}= src$ \\
\texttt{imul}      & \Rc/\Mc/\Lc & \Rc     & \Lc
  & Signed multiply (1/2/3-op) \\
\texttt{mul\{S\}}  & \Rc/\Mc & --- & ---
  & Unsigned mul: (R)DX:(R)AX = (R)AX $\times$ src \\
\texttt{div\{S\}}  & \Rc/\Mc & --- & ---
  & Unsigned div: (R)DX:(R)AX / src $\to$ RAX=quot, RDX=rem \\
\texttt{idiv\{S\}} & \Rc/\Mc & --- & ---
  & Signed div: (R)DX:(R)AX / src $\to$ RAX=quot, RDX=rem (use CDQ/CQO first) \\
\texttt{inc\{S\}}  & \Rc/\Mc & --- & --- & Increment by 1 \\
\texttt{dec\{S\}}  & \Rc/\Mc & --- & --- & Decrement by 1 \\
\texttt{neg\{S\}}  & \Rc/\Mc & --- & --- & Two's complement negate \\
\midrule
\texttt{and\{S\}}  & \Rc/\Mc/\Lc & \Rc/\Mc & --- & Bitwise AND \\
\texttt{or\{S\}}   & \Rc/\Mc/\Lc & \Rc/\Mc & --- & Bitwise OR \\
\texttt{xor\{S\}}  & \Rc/\Mc/\Lc & \Rc/\Mc & --- & Bitwise XOR \\
\texttt{not\{S\}}  & \Rc/\Mc & --- & --- & Bitwise NOT (ones' complement) \\
\midrule
\texttt{shl\{S\}}  & \Lc/\Rc(CL) & \Rc/\Mc & --- & Shift left \\
\texttt{shr\{S\}}  & \Lc/\Rc(CL) & \Rc/\Mc & --- & Shift right (logical) \\
\texttt{sar\{S\}}  & \Lc/\Rc(CL) & \Rc/\Mc & --- & Shift right (arithmetic) \\
\midrule
\texttt{cmp\{S\}}  & \Rc/\Mc/\Lc & \Rc/\Mc & ---
  & Compare: set flags on $dst - src$ \\
\texttt{test\{S\}} & \Rc/\Mc/\Lc & \Rc/\Mc & ---
  & Test: set flags on $dst \mathbin{\&} src$ \\
\midrule
\texttt{push\{S\}} & \Rc/\Mc/\Lc & --- & --- & Push onto stack (RSP $-=$ 8) \\
\texttt{pop\{S\}}  & \Rc/\Mc     & --- & --- & Pop from stack  (RSP $+=$ 8) \\
\midrule
\texttt{jmp}       & \Lc/\Mc/\Rc & --- & --- & Unconditional jump \\
\texttt{jcc}       & \Lc         & --- & --- & Conditional jump (see below) \\
\texttt{call}      & \Lc/\Mc/\Rc & --- & --- & Call subroutine \\
\texttt{ret}       & ---         & --- & --- & Return from subroutine \\
\midrule
\texttt{cbw}       & --- & --- & --- & Sign-extend AL $\to$ AX \\
\texttt{cwde}      & --- & --- & --- & Sign-extend AX $\to$ EAX \\
\texttt{cdqe}      & --- & --- & --- & Sign-extend EAX $\to$ RAX \\
\texttt{cwd}       & --- & --- & --- & Sign-extend AX $\to$ DX:AX \\
\texttt{cdq}       & --- & --- & --- & Sign-extend EAX $\to$ EDX:EAX \\
\texttt{cqo}       & --- & --- & --- & Sign-extend RAX $\to$ RDX:RAX \\
\midrule
\texttt{syscall}   & --- & --- & --- & Linux system call (RAX = number) \\
\texttt{nop}       & --- & --- & --- & No operation \\
\bottomrule
\end{tabular}
}

\columnbreak

% ---------- Intel Syntax ----------
{\small\bfseries\sffamily Intel Syntax}\par\vspace{-3pt}
{
\begin{tabular}{@{}l l l l p{4.2cm}@{}}
\toprule
\textbf{Mnem.} & \textbf{Op1} & \textbf{Op2} & \textbf{Op3}
  & \textbf{Description} \\
\midrule
\texttt{mov}       & \Rc/\Mc & \Rc/\Mc/\Lc & --- & Copy src$\to$dst \\
\texttt{movsx}     & \Rc     & \Rc/\Mc     & --- & Move with sign-extension \\
\texttt{movzx}     & \Rc     & \Rc/\Mc     & --- & Move with zero-extension \\
\texttt{lea}       & \Rc     & \Mc         & --- & Load effective address \\
\texttt{xchg}      & \Rc/\Mc & \Rc/\Mc     & --- & Exchange two operands \\
\midrule
\texttt{add}       & \Rc/\Mc & \Rc/\Mc/\Lc & ---
  & $dst \mathrel{+}= src$ \\
\texttt{sub}       & \Rc/\Mc & \Rc/\Mc/\Lc & ---
  & $dst \mathrel{-}= src$ \\
\texttt{imul}      & \Rc     & \Rc/\Mc     & \Lc
  & Signed multiply (1/2/3-op) \\
\texttt{mul}       & \Rc/\Mc & --- & ---
  & Unsigned mul: (R)DX:(R)AX = (R)AX $\times$ src \\
\texttt{div}       & \Rc/\Mc & --- & ---
  & Unsigned div: (R)DX:(R)AX / src $\to$ RAX=quot, RDX=rem \\
\texttt{idiv}      & \Rc/\Mc & --- & ---
  & Signed div: (R)DX:(R)AX / src $\to$ RAX=quot, RDX=rem (use CDQ/CQO first) \\
\texttt{inc}       & \Rc/\Mc & --- & --- & Increment by 1 \\
\texttt{dec}       & \Rc/\Mc & --- & --- & Decrement by 1 \\
\texttt{neg}       & \Rc/\Mc & --- & --- & Two's complement negate \\
\midrule
\texttt{and}       & \Rc/\Mc & \Rc/\Mc/\Lc & --- & Bitwise AND \\
\texttt{or}        & \Rc/\Mc & \Rc/\Mc/\Lc & --- & Bitwise OR \\
\texttt{xor}       & \Rc/\Mc & \Rc/\Mc/\Lc & --- & Bitwise XOR \\
\texttt{not}       & \Rc/\Mc & --- & --- & Bitwise NOT (ones' complement) \\
\midrule
\texttt{shl}       & \Rc/\Mc & \Lc/\Rc(CL) & --- & Shift left \\
\texttt{shr}       & \Rc/\Mc & \Lc/\Rc(CL) & --- & Shift right (logical) \\
\texttt{sar}       & \Rc/\Mc & \Lc/\Rc(CL) & --- & Shift right (arithmetic) \\
\midrule
\texttt{cmp}       & \Rc/\Mc & \Rc/\Mc/\Lc & ---
  & Compare: set flags on $dst - src$ \\
\texttt{test}      & \Rc/\Mc & \Rc/\Mc/\Lc & ---
  & Test: set flags on $dst \mathbin{\&} src$ \\
\midrule
\texttt{push}      & \Rc/\Mc/\Lc & --- & --- & Push onto stack (RSP $-=$ 8) \\
\texttt{pop}       & \Rc/\Mc     & --- & --- & Pop from stack  (RSP $+=$ 8) \\
\midrule
\texttt{jmp}       & \Lc/\Mc/\Rc & --- & --- & Unconditional jump \\
\texttt{jcc}       & \Lc         & --- & --- & Conditional jump (see below) \\
\texttt{call}      & \Lc/\Mc/\Rc & --- & --- & Call subroutine \\
\texttt{ret}       & ---         & --- & --- & Return from subroutine \\
\midrule
\texttt{cbw}       & --- & --- & --- & Sign-extend AL $\to$ AX \\
\texttt{cwde}      & --- & --- & --- & Sign-extend AX $\to$ EAX \\
\texttt{cdqe}      & --- & --- & --- & Sign-extend EAX $\to$ RAX \\
\texttt{cwd}       & --- & --- & --- & Sign-extend AX $\to$ DX:AX \\
\texttt{cdq}       & --- & --- & --- & Sign-extend EAX $\to$ EDX:EAX \\
\texttt{cqo}       & --- & --- & --- & Sign-extend RAX $\to$ RDX:RAX \\
\midrule
\texttt{syscall}   & --- & --- & --- & Linux system call (RAX = number) \\
\texttt{nop}       & --- & --- & --- & No operation \\
\bottomrule
\end{tabular}
}

\end{multicols}

\vspace{-2pt}
{\footnotesize
\textbf{Conditional jump (\texttt{jcc}) variants:}\quad
\texttt{je}/\texttt{jz}~($=$)\quad
\texttt{jne}/\texttt{jnz}~($\neq$)\quad
\texttt{jg}~(signed$>$)\quad
\texttt{jl}~(signed$<$)\quad
\texttt{jge}~($\geq$)\quad
\texttt{jle}~($\leq$)\quad
\texttt{ja}~(unsigned$>$)\quad
\texttt{jb}~(unsigned$<$)\quad
\texttt{jae}~($\geq$)\quad
\texttt{jbe}~($\leq$)\quad
\texttt{js}~(SF=1)\quad
\texttt{jns}~(SF=0)}

\end{tcolorbox}

\newpage

% ===================== PART 2 =====================
\begin{tcolorbox}[mysec,
  title={Part 2: General Purpose Registers \& Calling Conventions --- x86-64}]

{
\footnotesize
\begin{tabular}{@{}l l l p{2.9cm} p{2.9cm} p{2.9cm} p{2.5cm}@{}}
\toprule
\textbf{Reg} & \textbf{Alias} & \textbf{Saved}
  & \textbf{Role (General)}
  & \textbf{Microsoft x64}
  & \textbf{System V AMD64}
  & \textbf{Notes} \\
\midrule
RAX & accumulator & No
  & Return value; mul/div implicit
  & Return value (64-bit)
  & Return value (64-bit)
  & Caller-saved \\
RBX & base        & Yes
  & General purpose; callee-saved
  & Callee-saved
  & Callee-saved
  & Must preserve \\
RCX & counter     & No
  & Shift count (\texttt{cl}); loop
  & \textbf{Arg 1} (integer)
  & Arg 4 (integer)
  & Caller-saved \\
RDX & data        & No
  & I/O; mul/div high half
  & \textbf{Arg 2} (integer)
  & Arg 3 (integer)
  & Caller-saved \\
RSI & source idx  & No
  & String/memops source
  & Callee-saved (non-arg)
  & \textbf{Arg 2} (integer)
  & SysV arg; Win saved \\
RDI & dest idx    & No
  & String/memops destination
  & Callee-saved (non-arg)
  & \textbf{Arg 1} (integer)
  & SysV arg; Win saved \\
RSP & stack ptr   & ---
  & Top of stack (grows down)
  & Stack pointer
  & Stack pointer
  & 16-byte aligned \\
RBP & frame ptr   & Yes
  & Stack frame base (optional)
  & Callee-saved
  & Callee-saved
  & Often used as frame \\
\midrule
R8  & ---         & No
  & General purpose
  & \textbf{Arg 3} (integer)
  & \textbf{Arg 5} (integer)
  & Caller-saved \\
R9  & ---         & No
  & General purpose
  & \textbf{Arg 4} (integer)
  & \textbf{Arg 6} (integer)
  & Caller-saved \\
R10 & ---         & No
  & General purpose; static chain
  & Caller-saved; static chain ptr
  & Caller-saved
  & Win: used by closures \\
R11 & ---         & No
  & General purpose; scratch
  & Caller-saved; tail-call temp
  & Caller-saved
  & Volatile \\
R12 & ---         & Yes
  & General purpose
  & Callee-saved
  & Callee-saved
  & Must preserve \\
R13 & ---         & Yes
  & General purpose
  & Callee-saved
  & Callee-saved
  & Must preserve \\
R14 & ---         & Yes
  & General purpose
  & Callee-saved
  & Callee-saved
  & Must preserve \\
R15 & ---         & Yes
  & General purpose; TLS base
  & Callee-saved; TEB pointer
  & Callee-saved
  & Must preserve \\
\midrule
XMM0--5 & ---     & No
  & SSE 128-bit vector / FP args
  & \textbf{Arg 1--6} (float)
  & \textbf{Arg 1--8} (float)
  & Caller-saved \\
XMM6--15 & ---    & Yes / No
  & SSE 128-bit vector
  & \textbf{Callee-saved} (XMM6--15)
  & Caller-saved
  & Win: low 64-bit preserved \\
\midrule
RFLAGS & flags    & ---
  & CF ZF SF OF PF AF (status)
  & Volatile (except DF)
  & Volatile
  & Direction flag must be 0 \\
RIP & instr ptr   & ---
  & Address of next instruction
  & Not directly accessible
  & Not directly accessible
  & RIP-relative addressing \\
\bottomrule
\end{tabular}
}

\vspace{2pt}

\begin{minipage}[t]{0.48\textwidth}
\footnotesize\bfseries Microsoft x64 Calling Convention\mdseries\par\vspace{2pt}
\begin{tabular}{@{}l l@{}}
Integer args: & RCX, RDX, R8, R9 (first 4) \\
Float args:   & XMM0, XMM1, XMM2, XMM3 \\
Stack args:   & Pushed right-to-left starting at [RSP+32] \\
Shadow space: & \textbf{32 bytes} (caller must reserve) \\
Return (int): & RAX (64-bit) \\
Return (FP):  & XMM0 \\
Return (128): & RDX:RAX \\
Stack align:  & 16-byte before \texttt{call} \\
Callee-saved: & RBX, RBP, RDI, RSI, R12--R15, XMM6--XMM15 \\
\end{tabular}
\end{minipage}
\hfill
\begin{minipage}[t]{0.48\textwidth}
\footnotesize\bfseries System V AMD64 Calling Convention\mdseries\par\vspace{2pt}
\begin{tabular}{@{}l l@{}}
Integer args: & RDI, RSI, RDX, RCX, R8, R9 (first 6) \\
Float args:   & XMM0--XMM7 (first 8) \\
Stack args:   & Pushed right-to-left after [RSP+8] \\
Shadow space: & \textbf{None} (no reserved area) \\
Return (int): & RAX (64-bit) \\
Return (FP):  & XMM0 \\
Return (128): & RDX:RAX \\
Stack align:  & 16-byte before \texttt{call} \\
Callee-saved: & RBX, RBP, R12--R15 \\
\end{tabular}
\end{minipage}

\vspace{3pt}
\begin{minipage}[t]{0.48\textwidth}
\footnotesize
\textbf{Stack layout (Win x64):}\par
[RSP] = return addr\par
[RSP+8..31] = shadow space (32B, caller reserves for callee to spill regs)\par
[RSP+32..] = 5th+ args pushed right-to-left.\par
RSP must be 16-byte aligned before \texttt{call}.
\end{minipage}
\hfill
\begin{minipage}[t]{0.48\textwidth}
\footnotesize
\textbf{Stack layout (SysV AMD64):}\par
[RSP] = return addr\par
[RSP+8..] = 7th+ args pushed right-to-left.\par
No shadow space. No reserved area.\par
RSP must be 16-byte aligned before \texttt{call}.
\end{minipage}

\end{tcolorbox}

\newpage

% ===================== PART 3 =====================
\begin{tcolorbox}[mysec,
  title={Part 3: Program Templates --- Intel/Windows API \quad$|$\quad AT\&T/Linux API}]

\begin{multicols}{2}
\setlength{\parskip}{0pt}

% ---------- A: Win x64 API Template ----------
{\small\bfseries\sffamily A. Win x64 API Template (Intel Syntax)}\par
\vspace{1pt}
\begin{lstlisting}[style=asmcode]
# Build: gcc -no-pie win.s -o win.exe -lkernel32
#-----------------------------------------------
    .intel_syntax noprefix
    .section .data
msg:    .asciz "Hello!\n"
    .section .text
    .extern printf
    .extern exit
    .globl main
main:
    push    rbp
    mov     rbp, rsp
    sub     rsp, 32             # shadow space (32B)
    lea     rcx, msg[rip]       # arg1 (Win x64)
    xor     eax, eax
    call    printf
    xor     ecx, ecx            # arg1: exit code
    call    exit
\end{lstlisting}

% ---------- B: Win x64 ABI Function Demo ----------
{\small\bfseries\sffamily B. Win x64 Function Calling (Intel)}\par
\vspace{1pt}
\begin{lstlisting}[style=asmcode]
# --- Callee: receives args in RCX,RDX,R8 ---
sum3:
    push    rbp
    mov     rbp, rsp
    mov     rax, rcx            # rax = x
    add     rax, rdx            # rax += y
    add     rax, r8             # rax += z
    pop     rbp
    ret
# --- Caller: set args + shadow space ---
main:
    push    rbp
    mov     rbp, rsp
    sub     rsp, 32             # shadow space (32B)
    mov     ecx, 10             # arg1
    mov     edx, 20             # arg2
    mov     r8d, 30             # arg3
    call    sum3                # result in RAX
    add     rsp, 32
    pop     rbp
    ret
\end{lstlisting}

\columnbreak

% ---------- C: Linux API Template ----------
{\small\bfseries\sffamily C. Linux API Template (AT\&T Syntax)}\par
\vspace{1pt}
\begin{lstlisting}[style=asmcode]
# Build: gcc -no-pie lin.s -o lin
#-----------------------------------------------
    .section .data
msg:    .asciz "Hello!\n"
    .section .text
    .extern printf
    .extern exit
    .globl main
main:
    pushq   %rbp
    movq    %rsp, %rbp
    leaq    msg(%rip), %rdi     # arg1 (SysV)
    xorl    %eax, %eax
    call    printf
    xorl    %edi, %edi          # arg1: exit code
    call    exit
# Standalone (no libc):
# _start: movq $1,%rax; movq $1,%rdi
#   leaq msg(%rip),%rsi; movq $7,%rdx
#   syscall; movq $60,%rax; xorq %rdi,%rdi; syscall
\end{lstlisting}

% ---------- D: Linux ABI Function Demo ----------
{\small\bfseries\sffamily D. SysV Function Calling (AT\&T)}\par
\vspace{1pt}
\begin{lstlisting}[style=asmcode]
# --- Callee: receives args in RDI,RSI,RDX ---
sum3:
    pushq   %rbp
    movq    %rsp, %rbp
    movq    %rdi, %rax          # rax = x
    addq    %rsi, %rax          # rax += y
    addq    %rdx, %rax          # rax += z
    popq    %rbp
    ret
# --- Caller: set args, no shadow space ---
main:
    pushq   %rbp
    movq    %rsp, %rbp
    movl    $10, %edi           # arg1
    movl    $20, %esi           # arg2
    movl    $30, %edx           # arg3
    call    sum3                # result in RAX
    leave
    ret
\end{lstlisting}

\end{multicols}

\vspace{-1pt}
{\footnotesize
\textbf{Win x64:}\quad
RCX, RDX, R8, R9 (first 4 int args). Caller reserves 32B shadow space. 5th arg at [RSP+32].
\qquad
\textbf{SysV:}\quad
RDI, RSI, RDX, RCX, R8, R9 (first 6 int args). No shadow space. 7th+ arg on stack.}

\end{tcolorbox}

\newpage

% ===================== PART 4 =====================
\begin{tcolorbox}[mysec,
  title={Part 4: GAS Directives \& Pseudo-Operations}]

\begin{multicols}{2}
\setlength{\parskip}{0pt}

% ---------- Left: Section & Data directives ----------
{\small\bfseries\sffamily Section \& Data Directives}\par\vspace{-2pt}
{
\footnotesize
\begin{tabular}{@{}l p{5cm} p{4.5cm}@{}}
\toprule
\textbf{Directive} & \textbf{Description} & \textbf{Example} \\
\midrule
\texttt{.data}    & Begin initialized data segment
  & \texttt{.data} \\
\texttt{.section .data} & Same, explicit section name
  & \texttt{.section .data} \\
\texttt{.text}    & Begin code (text) segment
  & \texttt{.text} \\
\texttt{.section .text} & Same, explicit section name
  & \texttt{.section .text} \\
\texttt{.bss}     & Begin uninitialized data segment
  & \texttt{.bss} \\
\texttt{.section .bss} & Same, explicit section name
  & \texttt{.section .bss} \\
\midrule
\texttt{.byte}    & Emit 8-bit integer(s)
  & \texttt{.byte 0x41, 65} \\
\texttt{.word} / \texttt{.value} & Emit 16-bit integer(s)
  & \texttt{.word 1024} \\
\texttt{.long} / \texttt{.int} & Emit 32-bit integer(s)
  & \texttt{.long 0xDEADBEEF} \\
\texttt{.quad}    & Emit 64-bit integer(s)
  & \texttt{.quad 0x123456789} \\
\texttt{.ascii}   & Emit string (no null terminator)
  & \texttt{.ascii "Hello"} \\
\texttt{.asciz} / \texttt{.string} & Emit string with null terminator
  & \texttt{.asciz "Hello"} \\
\texttt{.float}   & Emit 32-bit IEEE float
  & \texttt{.float 3.14} \\
\texttt{.double}  & Emit 64-bit IEEE double
  & \texttt{.double 2.71828} \\
\texttt{.space} / \texttt{.skip} & Emit $n$ zero bytes
  & \texttt{.space 256} \\
\texttt{.fill}    & Fill $n$ copies of $s$-byte value $v$
  & \texttt{.fill 10, 4, 0} \\
\midrule
\texttt{.zero}    & Zero-fill in BSS (no file size)
  & \texttt{.zero 4096} \\
\bottomrule
\end{tabular}
}

\columnbreak

% ---------- Right: Symbol & Control directives ----------
{\small\bfseries\sffamily Symbol \& Control Directives}\par\vspace{-2pt}
{
\footnotesize
\begin{tabular}{@{}l p{5cm} p{4.5cm}@{}}
\toprule
\textbf{Directive} & \textbf{Description} & \textbf{Example} \\
\midrule
\texttt{.globl} / \texttt{.global} & Export symbol (visible to linker)
  & \texttt{.globl main} \\
\texttt{.extern}  & Declare external symbol (optional in GAS)
  & \texttt{.extern printf} \\
\midrule
\texttt{.equ}     & Define constant (immutable)
  & \texttt{.equ BUFSIZ, 1024} \\
\texttt{.set}     & Define symbol (mutable)
  & \texttt{.set counter, 0} \\
\texttt{.equiv}   & Define constant (error if redefined)
  & \texttt{.equiv TRUE, 1} \\
\midrule
\texttt{.align}   & Align to $2^n$ boundary
  & \texttt{.align 4} (=16 bytes) \\
\texttt{.balign}  & Align to $n$-byte boundary
  & \texttt{.balign 16} \\
\texttt{.p2align} & Align to $2^n$ boundary (explicit)
  & \texttt{.p2align 4} \\
\midrule
\texttt{.type}    & Set ELF symbol type
  & \texttt{.type main, @function} \\
\texttt{.size}    & Set ELF symbol size
  & \texttt{.size main, .-main} \\
\texttt{.section} & Switch to named section with flags
  & \texttt{.section .rodata,"a"} \\
\midrule
\texttt{.intel\_syntax noprefix} & Switch to Intel syntax (GAS)
  & \texttt{.intel\_syntax noprefix} \\
\texttt{.att\_syntax prefix} & Switch back to AT\&T syntax (GAS)
  & \texttt{.att\_syntax prefix} \\
\midrule
\texttt{.org}     & Set location counter
  & \texttt{.org 0x100} \\
\texttt{.incbin}  & Include raw binary file
  & \texttt{.incbin "data.bin"} \\
\texttt{.macro} / \texttt{.endm} & Define macro
  & \texttt{.macro pusha} ... \texttt{.endm} \\
\texttt{.rept} / \texttt{.endr} & Repeat block $n$ times
  & \texttt{.rept 16} ... \texttt{.endr} \\
\midrule
\texttt{.code16} / \texttt{.code32} / \texttt{.code64} & Set instruction mode
  & \texttt{.code64} \\
\texttt{.file}    & Set source file name (debug info)
  & \texttt{.file "main.s"} \\
\texttt{.loc}     & Emit debug line number
  & \texttt{.loc 1 42} \\
\bottomrule
\end{tabular}
}

\end{multicols}

% ---------- Macro Demo ----------
\vspace{2pt}
{\small\bfseries\sffamily Macro Demonstration (AT\&T Syntax)}\par
\vspace{1pt}
\begin{multicols}{2}
\setlength{\parskip}{0pt}
\begin{lstlisting}[style=asmcode]
# Simple: print via Linux syscall
.macro print_str str, len
    movq  $1, %rax
    movq  $1, %rdi
    leaq  \str(%rip), %rsi
    movq  $\len, %rdx
    syscall
.endm

# Default param: exit(code=0)
.macro exit_prog code=0
    movq  $60, %rax
    movq  $\code, %rdi
    syscall
.endm
\end{lstlisting}

\begin{lstlisting}[style=asmcode]
# Local labels: .L + \@ unique ID
.macro safe_div dividend, divisor
    cmpq  $0, \divisor
    je    .Ldiv_zero\@
    movq  \dividend, %rax
    cqo
    idivq \divisor
    jmp   .Ldiv_done\@
.Ldiv_zero\@:
    xorq  %rax, %rax
.Ldiv_done\@:
.endm

# Usage:
_start:
    print_str msg, 3
    safe_div %rbx, %rcx
    exit_prog 0
\end{lstlisting}
\end{multicols}

\vspace{2pt}
{\footnotesize
\textbf{Notes:}
\textbullet\enspace In GAS, \texttt{.data}, \texttt{.text}, \texttt{.bss} can be used as shorthand or with \texttt{.section}.
\textbullet\enspace \texttt{.extern} is optional in GAS (all undefined symbols are implicitly external), but improves readability.
\textbullet\enspace \texttt{.align $n$} aligns to $2^n$ bytes on most targets; use \texttt{.balign $n$} for unambiguous byte alignment.
\textbullet\enspace Use \texttt{.intel\_syntax noprefix} at the top of a GAS file to switch to Intel syntax; default is AT\&T.
\textbullet\enspace \texttt{.type}, \texttt{.size}, \texttt{.section} flags (e.g.\ \texttt{"a"@progbits, "w"@nobits}) are ELF-specific; Windows/PE uses \texttt{.def}/\texttt{.scl}/\texttt{.endef}.}

\end{tcolorbox}

% ===================== PART 5 =====================
\begin{tcolorbox}[mysec,
  title={Part 5: GCC / Clang Extended Inline Assembly --- x86-64}]

% ---------- Constraints Table ----------
{\small\bfseries\sffamily Operand Constraint Reference}\par
\vspace{-2pt}
{
\scriptsize
\begin{tabular}{@{}l p{4.5cm} l p{4.5cm} l p{4cm}@{} l@{}}
\toprule
\textbf{Cstr.} & \textbf{Meaning} & \textbf{Size}
  & \textbf{Cstr.} & \textbf{Meaning} & \textbf{Size} \\
\midrule
\multicolumn{6}{@{}l@{}}{\textit{Register class constraints}} \\
\texttt{r}  & Any general-purpose reg       & 1/2/4/8
  & \texttt{v}  & Any vector reg (SSE/AVX)     & 16/32 \\
\texttt{q}  & Byte-addressable reg (a,b,c,d) & 1/2/4/8
  & \texttt{x}  & Any SSE reg (XMM)            & 16 \\
  & \texttt{y}  & Any MMX register              & 8 \\
\midrule
\multicolumn{6}{@{}l@{}}{\textit{Specific register constraints}} \\
\texttt{a}  & RAX / EAX / AX / AL           & 1/2/4/8
  & \texttt{S}  & RSI / ESI / SI / SIL         & 1/2/4/8 \\
\texttt{b}  & RBX / EBX / BX / BL           & 1/2/4/8
  & \texttt{D}  & RDI / EDI / DI / DIL         & 1/2/4/8 \\
\texttt{c}  & RCX / ECX / CX / CL           & 1/2/4/8
  & \texttt{d}  & RDX / EDX / DX / DL          & 1/2/4/8 \\
\texttt{A}  & RDX:RAX pair (for \texttt{div}) & 16
  & ---   & ---                            & --- \\
\midrule
\multicolumn{6}{@{}l@{}}{\textit{Immediate \& memory constraints}} \\
\texttt{i}  & Integer immediate (compile-time) & ---
  & \texttt{m}  & Any memory operand             & --- \\
\texttt{n}  & Integer immediate (run-time)    & ---
  & \texttt{p}  & Valid memory address (pointer)  & 8 \\
\texttt{I}  & Immediate 0--31 (for shift)     & ---
  & \texttt{o}  & Offsetable memory operand       & --- \\
\texttt{J}  & Immediate 0--63                 & ---
  & \texttt{V}  & Non-offsetable memory           & --- \\
\texttt{K}  & Immediate $-128$--$127$ (\texttt{movabs}) & ---
  & \texttt{X}  & Any operand                     & --- \\
\texttt{L}  & Immediate $\pm 2^{31}$ (for \texttt{jmp}) & ---
  & \texttt{g}  & Any register or memory          & --- \\
\texttt{M}  & Immediate 0--3 (for \texttt{lea}) & ---
  & \texttt{E}  & Float immediate (standard)      & --- \\
\texttt{N}  & Immediate $-128$--$127$ (for \texttt{in}) & ---
  & \texttt{F}  & Float immediate (standard)      & --- \\
\bottomrule
\end{tabular}
}

\vspace{2pt}

% ---------- Modifiers Table ----------
{\small\bfseries\sffamily Output Modifiers (x86)}\par
\vspace{-2pt}
{
\scriptsize
\begin{tabular}{@{}l p{6cm} l p{6cm}@{} l@{}}
\toprule
\textbf{Mod.} & \textbf{Effect (on output operand)}
  & \textbf{Mod.} & \textbf{Effect (on output operand)} \\
\midrule
\texttt{c}  & Print constant without \texttt{\$} prefix
  & \texttt{P}  & Print only the raw constant (no \texttt{\$}, no \texttt{\%}) \\
\texttt{h}  & Print high byte of register (e.g.\ \texttt{ah})
  & \texttt{[0-9]} & Print digit \texttt{N} of the operand (sub-operand) \\
\texttt{b}  & Print low byte register name (8-bit)
  & \texttt{w}  & Print 16-bit register name \\
\texttt{k}  & Print 32-bit register name
  & \texttt{q}  & Print 64-bit register name \\
\bottomrule
\end{tabular}
}

\vspace{2pt}

% ---------- Inline ASM Syntax ----------
{\small\bfseries\sffamily Syntax: \texttt{\_\_asm\_\_\_ ( "template" : output : input : clobbers )}}\par
\vspace{-2pt}
{
\scriptsize
\begin{tabular}{@{}l p{10cm}@{} l@{}}
\toprule
\textbf{Field} & \textbf{Description} \\
\midrule
\texttt{template}  & Assembly instructions with \texttt{\%0}, \texttt{\%1}, ...\ operand placeholders. Use \texttt{\%k0} for 32-bit name, \texttt{\%q0} for 64-bit. \\
\texttt{output}    & \texttt{"=r"(var)} write-only, \texttt{"=rm"(var)} reg-or-mem. \texttt{"+r"(var)} read-write. \\
\texttt{input}     & \texttt{"r"(val)} register, \texttt{"i"(val)} immediate, \texttt{"m"(val)} memory, \texttt{"0"(var)} bind to output 0. \\
\texttt{clobbers}  & \texttt{"memory"} (writes to memory), \texttt{"cc"} (modifies flags), \texttt{"rax"} (trashes reg). \\
\bottomrule
\end{tabular}
}

\vspace{2pt}

% ---------- EFLAGS Example ----------
{\small\bfseries\sffamily Example: Capture EFLAGS to a C variable}\par
\vspace{-2pt}
\begin{multicols}{2}
\setlength{\parskip}{0pt}
\begin{lstlisting}[style=asmcode, basicstyle=\ttfamily\scriptsize]
#include <stdint.h>

// Method 1: lahf
// captures SF,ZF,AF,PF,CF into AH
static inline uint8_t
get_flags_lahf(void) {
    uint8_t f;
    __asm__ volatile ("lahf"
        : "=a"(f) : : "memory");
    return f; // bits: CF=0 PF=2
              //        AF=4 ZF=6 SF=7
}

// Method 2: pushfq/popq
// captures full 64-bit RFLAGS
static inline uint64_t
get_rflags(void) {
    uint64_t f;
    __asm__ volatile (
        "pushfq\n\tpopq %0"
        : "=r"(f) : : "memory");
    return f;
}
\end{lstlisting}

\begin{lstlisting}[style=asmcode, basicstyle=\ttfamily\scriptsize]
// Usage: test if a == b
int test_eq(int a, int b) {
    uint64_t flags;
    __asm__ (
        "cmpl %2, %1\n\t"
        "pushfq\n\tpopq %0"
        : "=r"(flags)
        : "r"(a), "r"(b)
        : "cc"
    );
    // ZF bit is bit 6
    return (flags & 0x40) != 0;
}
\end{lstlisting}
\end{multicols}

\vspace{2pt}
{\scriptsize
\textbf{Key points:}
\textbullet\enspace \texttt{\%0}, \texttt{\%1}... reference operands in GAS (AT\&T) syntax by default; add \texttt{\%} as prefix for Intel via \texttt{\%[name]}.
\textbullet\enspace Use \texttt{\%k0} to force 32-bit register name (e.g.\ \texttt{eax}), \texttt{\%q0} for 64-bit (e.g.\ \texttt{rax}).
\textbullet\enspace \texttt{volatile} prevents compiler from removing or reordering the asm block.
\textbullet\enspace \texttt{"cc"} clobber tells the compiler that flags are modified.
\textbullet\enspace \texttt{lahf} is fast but only captures 5 flags; \texttt{pushfq}/\texttt{popq} captures all 64 bits of RFLAGS.}

\end{tcolorbox}

\end{document}
